\documentclass[11pt]{letter}
\usepackage[T1]{fontenc}
\usepackage[utf8]{inputenc}
\usepackage{lmodern}
\usepackage[margin=1in]{geometry}
\usepackage{microtype}

\signature{Darcio Genicolo-Martins\\INSPER, S\~ao Paulo\\\texttt{darciogm1@insper.edu.br}}
\address{Darcio Genicolo-Martins\\INSPER\\R.\ Quat\'a 300, S\~ao Paulo, SP, 04546-042, Brazil}

\begin{document}

\begin{letter}{The Editor\\\emph{Journal of Public Economics}}

\opening{Dear Editor,}

I am pleased to submit my manuscript, \emph{The Cost of Inclusion: Decomposing Bidder Exclusion in Public Procurement}, for consideration at the \emph{Journal of Public Economics}.

The paper studies the welfare consequences of small-and-medium-enterprise (SME) procurement set-asides, an instrument deployed in over a hundred jurisdictions worldwide and made a procurement default by Brazilian Law 14{,}133/2021. A decade of empirical evidence places the price cost of bidder-exclusion rules at 10--15 percent of winning auctions, but the mechanism behind the markup---and therefore the welfare ranking of bidder exclusion against alternative preference instruments---has remained out of reach because price regressions cannot separate the change in \emph{who} can bid from the change in \emph{how} each remaining bidder bids. The paper closes that gap on a clean total-exclusion natural experiment: a March 2018 administrative reversal that extended an existing SME-only rule to medical and hospital supplies on São Paulo's centralized R\$13 billion procurement platform, leaving 76 product groups under their pre-existing regimes as never-treated controls. Drop-out bids from the platform's Pregão English-reverse format point-identify type-specific cost distributions under the Haile--Tamer interpretation; an asymmetric IPV counterfactual conducted under observed equilibrium entry decomposes the simulated price effect into a within-auction component and a participation-margin component on the same footing where the partial-preference structural literature (Marion 2007; Krasnokutskaya 2011; Athey, Levin, and Seira 2011) leaves the two margins entangled.

Three results organize the paper. First, the simulated set-aside price effect operates predominantly within the auction under the model's main specification, with the within-auction component (which I label \emph{sheltered bidding}) accounting for 73--85 percent of the simulated total effect across cost-distribution estimators and across two specifications of the post-policy SME pool. Second, endogenous SME entry is the government's partial offset rather than the source of the markup: a fixed-pool counterfactual would deliver a simulated effect roughly 50--60 percent larger than the realized one, which means a fixed-pool welfare calculation---the standard reduced-form shortcut---materially misstates the realized policy cost. Third, a 10 percent price preference welfare-dominates the set-aside in thick, standardized non-pharmaceutical markets across both the main and the strict-invariance specifications, at near-zero fiscal cost; in thin, heterogeneous pharmaceutical markets the ranking is conditional on how the post-policy SME pool is modeled. The conditional reading is itself a policy-design statement: bidder-exclusion rules are most model-sensitive precisely where the protected pool is thin and the good is heterogeneous, and that is where caution is needed in deploying them.

I see the paper as a public-economics contribution, distinct from the structural empirical-IO literature it dialogues with. The exercise translates a procurement-policy reform into a welfare-decomposition object on which standard public-finance arithmetic (allocative deadweight loss, marginal cost of public funds) can be applied; the policy comparison against a 10 percent price preference is the kind of instrument-versus-instrument question that JPubE has historically sponsored. The Brazilian setting is a plausibly exogenous total-exclusion design, providing the data structure that the U.S.-federal and timber settings of the prior structural procurement-preference literature could not. The paper sits in dialogue with recent JPubE-adjacent work on emerging-market procurement (Decarolis et al. 2024; Cabral 2017; Szerman 2023) and on procurement frictions and capacity-building (Bosio et al. 2022).

A difference-in-differences against 76 never-treated product groups, reported as a design-validation appendix in the manuscript, fixes the magnitude at 10--11 percent on absolute winning prices; the structural counterfactual on its structural normalization is roughly three to four times that DiD coefficient ($\sim$34--47 percent of the open-regime baseline $p_{S_1}$), reflecting that the two objects average over distinct primitives (realized log final prices with item and month fixed effects vs.\ the BNE-equivalent transaction price on UH-cleaned cost residuals). The structural exercise of this paper decomposes that effect into within-auction and participation-margin channels and prices the welfare cost (28.7 percent of $p_{S_1}$ in non-pharma and 47.0 percent in pharma at $\lambda = 0.30$). Scaled to the post-policy run-rate of Group-65 procurement, the implied welfare loss is on the order of R\$\,55--128 million (US\$\,16--37 million) per year on the medical-supplies group alone---an aggregate burden that the structural decomposition shows to be over 70 percent attributable to within-auction sheltered bidding, not to entry margin reshuffling.

The manuscript is approximately 50 pages of body text plus references, an online appendix, and supplementary tables. Replication code, anonymized analytical sample, and data dictionary will be made available at the project repository upon acceptance; the underlying BEC microdata are available from SEFAZ-SP under a data-use agreement.

The paper is not under review at any other journal and has not been previously published in any form. I have no conflicts of interest to declare.

I look forward to your editorial decision.

\closing{Sincerely,}

\end{letter}

\end{document}
